\documentclass[english]{article}
\usepackage[T1]{fontenc}
\usepackage[utf8]{inputenc}
\setcounter{secnumdepth}{-2}
\setlength{\parindent}{0bp}
\usepackage{amsmath}
\usepackage{amsthm}
\usepackage{amssymb}

\makeatletter
\theoremstyle{definition}
\newtheorem*{xca*}{\protect\exercisename}
\theoremstyle{plain}
\newtheorem*{thm*}{\protect\theoremname}
\theoremstyle{remark}
\newtheorem*{rem*}{\protect\remarkname}
\theoremstyle{plain}
\newtheorem{thm}{\protect\theoremname}
\theoremstyle{plain}
\newtheorem{lem}[thm]{\protect\lemmaname}

\@ifundefined{date}{}{\date{}}
\usepackage{graphicx}
\usepackage{geometry}
\newcommand{\limi}{\underset{n\rightarrow\infty}{\lim}}
\newcommand{\limxz}{\underset{x\rightarrow x_{0}}{\lim}}
\newcommand{\limfxz}{\underset{x\rightarrow x_{0}}{\lim}f(x)}
\newcommand{\limfxm}{\underset{x\rightarrow x_{0}^{-}}{\lim}f(x)}
\newcommand{\limfxp}{\underset{x\rightarrow x_{0}^{+}}{\lim}f(x)}
\newcommand{\limxm}{\underset{x\rightarrow x_{0}^{-}}{\lim}}
\newcommand{\limxp}{\underset{x\rightarrow x_{0}^{+}}{\lim}}
\newcommand{\sqrm}{M_{m\times m}(\mathbb{F})}
\newcommand{\seqa}{(a_{n})_{n=1}^{\infty}}
\newcommand{\seqx}{(x_{n})_{n=1}^{\infty}}
\newcommand{\funcdr}{f:D\rightarrow\mathbb{R}}
\newcommand{\der}{\limxz\frac{f(x)-f(x_{0})}{x-x_{0}}}
\usepackage[all]{pdfprivacy}

\makeatother

\usepackage{babel}
\providecommand{\exercisename}{Exercise}
\providecommand{\lemmaname}{Lemma}
\providecommand{\remarkname}{Remark}
\providecommand{\theoremname}{Theorem}

\begin{document}
\newgeometry{left=2cm, right=2cm, top=2cm, bottom=2cm}
\title{{\Large\vspace{-2em}Exercises 2B, Q1\vspace{-2em}}}
\maketitle
\begin{xca*}
Determine all vector spaces that have exactly one basis, even up to
reordering, and prove that no others exist. Note that this is different
from the original question. 
\end{xca*}
\begin{thm*}
A vector space $V$ over $\mathbb{F}$ has exactly one basis, even
up to reordering, if and only if it satisfies one of the two following
statements:

1. $V=\left\{ 0\right\} $ over any field $\mathbb{F}$

2. $\mathbb{F}$ is isomorphic to $\mathbb{Z}_{2}$ and $\exists0\ne v\text{ s.t. }V=\left\{ 0,v\right\} $.
\end{thm*}
\begin{rem*}
In the book, a basis is defined as a list, which is a finite object. 

Therefore, exactly one basis, even up to reordering, means that every
basis of the vector space is a permutation of a single list.
\end{rem*}
\begin{rem*}
If you did not learn about isomorphism and prime fields, you can think
of ``$\mathbb{F}$ is isomorphic to $\mathbb{Z}_{2}$'' as $\mathbb{F}=\{0,1\}$
with $1+1=0$.
\end{rem*}
\begin{lem}
If $V$ isn't finite-dimensional, then V has no basis. In particular,
V does not have exactly one basis, even up to reordering.
\end{lem}

\begin{proof}
According to definition $2.9$, a vector space is finite-dimensional
iff some list of vectors in it spans the space. 

Therefore, if it isn't finite-dimensional (and by definition $2.13$
it is then called infinite-dimensional), no list of vectors in it
spans the space.

Since by definition $2.26$ a basis must span the space, no list of
vectors can be a basis of the space, hence it has no bases and in
particular doesn't have exactly one. 
\end{proof}
\begin{lem}
$V=\left\{ 0\right\} $ over any field $\mathbb{F}$ has exactly one
basis.
\end{lem}

\begin{proof}
$\left\{ 0\right\} $ is a vector space over any field $\mathbb{F}$
- after definition $1.22$, $\left\{ 0\right\} $ is mentioned as
the simplest vector space.

Definition $2.4$ asserts $\text{Span }()=\left\{ 0\right\} $.

Definition $2.15$ asserts that the empty list is linearly independent. 

Therefore, by definition $2.26$, $\left(\right)$ is a basis for
$\left\{ 0\right\} $.

Any non-empty list of vectors in $\left\{ 0\right\} $ must contain
the zero vector, which makes it linearly dependent (according to example
$2.18$).

Thus, no non-empty list is a basis of $\left\{ 0\right\} $ (by definition
$2.26$), so $\left(\right)$ is its only basis.
\end{proof}
\newpage{}
\begin{lem}
Let $V$ be a finite-dimensional vector space over field $\mathbb{F}$,
and let $\mathcal{B}=\left(v_{1},\dots,v_{n}\right)$ be a basis of
$V$ such that $n=1$. 

Then $V$ has exactly one basis, even up to reordering, if and only
if $\mathbb{F}$ is isomorphic to $\mathbb{Z}_{2}$.
\end{lem}

\begin{proof}
If $\mathbb{F}$ is isomorphic to $\mathbb{Z}_{2}$, then $\mathbb{F}=\left\{ 0,1\right\} $
with $1+1=0$, and
\[
V=\text{Span }\mathcal{B}=\left\{ av_{1}\mid a\in\mathbb{F}\right\} =\left\{ 0,v_{1}\right\} 
\]
and there is no other basis other than $\mathcal{B}=\left(v_{1}\right)$:

Since $\mathcal{B}$ is linearly independent, $v_{1}\ne0$ (by example
$2.18$).

The empty list is not a basis, since by definition $2.4$, $\text{Span }\left(\right)=\left\{ 0\right\} \ne V$.

Every vector in $V$ is $0$ or $v_{1}$, so every list of vectors
in $V$ would contain some combination of these two elements.

By example $2.18$, any list containing $0$, or containing $v_{1}$
more than once, is linearly dependent and therefore not a basis.

Hence, overall, the only possible basis remains $\left(v_{1}\right)$.

If $\mathbb{F}$ is not isomorphic to $\mathbb{Z}_{2}$, then $\mathbb{F}$
has more than two elements, since every field with exactly two elements
is isomorphic to $\mathbb{Z}_{2}$. Hence, there exists $c\in\mathbb{F}$
with $c\notin\{0,1\}$.

Because $c\notin\left\{ 0,1\right\} $ and $v_{1}$ cannot be $0$
($\mathcal{B}$ would be linearly dependent), $cv_{1}\ne v_{1}$:

Otherwise, 
\[
0=v_{1}-cv_{1}\iff0=\left(1-c\right)v_{1}\underset{\text{Exercises 1B Q 2}}{\iff}1-c=0\iff c=1
\]
which is a contradiction.

Hence, $\left\{ cv_{1}\right\} \ne\left\{ v_{1}\right\} $, and we
will show that $\mathcal{B}'=\left(cv_{1}\right)$ is also a basis
for $V$ (we have shown that they are different even up to reordering). 

Because $v_{1}$ isn't the zero vector and $c\ne0$, the contra-positive
of exercises $1B$ question $2$ gives $cv_{1}\ne0$, and $\mathcal{B}'$
contains a single non-zero vector, hence it is linearly independent.

Additionally, 
\[
V=\text{Span }\mathcal{B}=\left\{ av_{1}\mid a\in\mathbb{F}\right\} =\left\{ a\cdot1\cdot v_{1}\mid a\in\mathbb{F}\right\} \underset{0\ne c\in\mathbb{F}}{=}\left\{ a\cdot\frac{1}{c}\cdot c\cdot v_{1}\mid a\in\mathbb{F}\right\} =\left\{ \frac{a}{c}\cdot cv_{1}\mid a\in\mathbb{F}\right\} \subseteq\text{Span }\mathcal{B}'
\]
and since $\text{Span }\mathcal{B}'$ is a subspace of $V$ (by $2.6$),
the two sided inclusion means $\mathcal{B}'$ spans $V$, and therefore
it is a basis for $V$.

This means that if $\mathbb{F}$ is not isomorphic to $\mathbb{Z}_{2}$,
$V$ has more than one basis, even up to reordering.
\end{proof}
\newpage{}
\begin{lem}
Let $V$ be a finite-dimensional vector space over field $\mathbb{F}$,
and let $\mathcal{B}=\left(v_{1},\dots,v_{n}\right)$ be a basis of
$V$ such that $n\ge2$. Then $\mathcal{B}$ is not the only basis
of $V$, even up to reordering. 
\end{lem}

\begin{proof}
Let us look at $\mathcal{B}'=\left(v_{1}+v_{2},v_{2},\dots,v_{n}\right)$
and let $v\in V$.

According to definition $2.26$, since $\mathcal{B}$ is a basis of
$V$, $\mathcal{B}$ spans $V$. 

Therefore, 
\[
\exists a_{1},a_{2},a_{3},\dots,a_{n}\in\mathbb{F}\ \text{ s.t. }v=a_{1}v_{1}+a_{2}v_{2}+a_{3}v_{3}+\dots+a_{n}v_{n}
\]

We can select
\[
\begin{cases}
c_{1}=a_{1} & \in\mathbb{F}\\
c_{2}=a_{2}-a_{1} & \in\mathbb{F}\\
c_{3}=a_{3} & \in\mathbb{F}\\
\vdots\\
c_{n}=a_{n} & \in\mathbb{F}
\end{cases}
\]

we get 
\[
v=a_{1}v_{1}+a_{2}v_{2}+a_{3}v_{3}+\dots+a_{n}v_{n}=a_{1}v_{1}+a_{1}v_{2}+a_{2}v_{2}-a_{1}v_{2}+a_{3}v_{3}+\dots+a_{n}v_{n}
\]
\[
=c_{1}\left(v_{1}+v_{2}\right)+c_{2}v_{2}+c_{3}v_{3}+\dots+c_{n}v_{n}\in\text{Span }\mathcal{B}'
\]
Hence $\mathcal{B}'$ also spans $V$.

Let $b_{1},b_{2},b_{3},\dots,b_{n}\in\mathbb{F}$ s.t. 
\[
b_{1}\left(v_{1}+v_{2}\right)+b_{2}v_{2}+b_{3}v_{3}+\dots+b_{n}v_{n}=0
\]

Then
\[
b_{1}v_{1}+\left(b_{1}+b_{2}\right)v_{2}+\dots+b_{n}v_{n}=0
\]

According to definition $2.26$, since $\mathcal{B}$ is a basis of
$V$, $\mathcal{B}$ is linearly independent and
\[
b_{1}=0,\quad\left(b_{1}+b_{2}\right)=0,\quad\dots,\quad b_{n}=0
\]
hence 
\[
b_{1}=b_{2}=\dots=b_{n}=0
\]

and according to definition $2.15$, $\mathcal{B}'$ is linearly independent
as well.

According to definition $2.26$, $\mathcal{B}'$ is a basis of $V$.

$\mathcal{B}'\ne\mathcal{B}$, nor is it a reordering of $\mathcal{B}$:
let us assume for the sake of contradiction that $v_{1}+v_{2}\in\mathcal{B}$.

Then $\exists k\in\left\{ 1,\dots,n\right\} \text{ s.t. }v_{k}=v_{1}+v_{2}$.

If $k=1$ or $k=2$ then accordingly $v_{1}+v_{2}=v_{1}$ or $v_{1}+v_{2}=v_{2}$
hence $v_{2}$ or $v_{1}$ is the zero vector, contradicting the linear
independence of $\mathcal{B}$ (according to example $2.18$, any
list containing the zero vector is linearly dependent).

If $k\ge3$ then $\mathcal{B}$ is of the form $\left(v_{1},v_{2},\dots,v_{1}+v_{2},\dots,v_{n}\right)$
and that 
\[
1\cdot v_{1}+1\cdot v_{2}+0\cdot v_{3}+\dots+0\cdot v_{k-1}+\left(-1\right)\cdot v_{k}+0\cdot v_{k+1}+\dots+0\cdot v_{n}=0
\]
contradicting the linear independence of $\mathcal{B}$.

Hence $\mathcal{B}'$ contains a vector not in $\mathcal{B}$, so
$\mathcal{B}'$ is a basis of $V$ different from $\mathcal{B}$,
even up to reordering.
\end{proof}
\newpage{}
\begin{proof}
\textbf{Tying it all together, }

Let $V$ be a vector space over field $\mathbb{F}$.

$\Leftarrow$: Let us assume $V$ satisfies one of the two following
statements:

1. $V=\left\{ 0\right\} $ over any field $\mathbb{F}$

2. $\mathbb{F}$ is isomorphic to $\mathbb{Z}_{2}$ and $\exists0\ne v\text{ s.t. }V=\left\{ 0,v\right\} $.

In the first case, we proved in lemma 2 that it has exactly one basis,
even up to reordering.

In the second case, $\left(v\right)$ is a basis of $V$ and we proved
in lemma 3 that $V$ has exactly one basis, even up to reordering.

$\Rightarrow\text{ via contra-positive}:$ Let us assume that $V$
satisfies neither of these statements, and cover all possible cases
of $V$.

If $V$ is not finite-dimensional, then we proved in lemma 1 that
it doesn't have exactly one basis.

Otherwise, $V$ is finite-dimensional, therefore it has a basis.

Let $\mathcal{B}=\left(v_{1},\dots,v_{n}\right)$ be a basis of $V$.

If $n=0$, then $V=\text{Span }\mathcal{B}=\text{Span }\left(\right)=\left\{ 0\right\} $
and it contradicts the fact that $V$ doesn't satisfy statement 1.

If $n=1$ and $\mathbb{F}$ isn't isomorphic to $\mathbb{Z}_{2}$,
then we proved in lemma 3 that $V$ doesn't have exactly one basis.

If $n=1$ and $\mathbb{F}$ is isomorphic to $\mathbb{Z}_{2}$, then
\[
V=\text{Span }\left(v_{1}\right)=\left\{ av_{1}\mid a\in\mathbb{F}\right\} =\left\{ 0\cdot v_{1},1\cdot v_{1}\right\} =\left\{ 0,v_{1}\right\} 
\]

and $v_{1}\ne0$ because $\mathcal{B}$ is a basis and is therefore
linearly independent - in contradiction to the fact that $V$ doesn't
satisfy statement 2.

If $n\ge2$, then we proved in lemma 4 that $V$ doesn't have exactly
one basis.

This covers all of the possibilities, as required. 
\end{proof}

\end{document}
